% -----------------------------------------------------------
\section{Love}
% -----------------------------------------------------------
	\label{LOVE}
	
% % ----------------------
% \subsection{See also}
% % ----------------------

% % ----------------------
% \subsection{1828 American Dictionary of the English Language} % *1828
% % ----------------------
	% \label{LOVE-1828}

	% \begin{compactitem}
		% \item
	% \end{compactitem}

% % ----------------------
% \subsection{Oxford English Dictionary} % *OED
% % ----------------------
	% \label{LOVE-OED}

	% \begin{compactitem}
		% \item[]
	% \end{compactitem}
	
% % ----------------------
% \subsection{Restoration Lexicon} % *RL *RESTORATION LEXICON
% % ----------------------
	% \label{LOVE-OED}

	% \begin{compactitem}
		% \item[]
	% \end{compactitem}

% % ----------------------
% \subsection{Scriptural Usage} % *SCRIPTURES
% % ----------------------
	% \label{LOVE-SCRIPTURES}

	% % -----
	% \subsubsection{Old Testament} % *OT
	% % -----
		% \label{LOVE-OT}
	
		% % --
		% \paragraph{KJV}
		% % --
			% \label{LOVE-OT-KJV}
	
			% \begin{tabular}{ll}
				% Total occurrences in the OT & 
			% \end{tabular}
			
			% \begin{compactdesc}
				% \item[]
			% \end{compactdesc}
			
		% % --
		% \paragraph{Hebrew Bible}
		% % --
			% \label{LOVE-HEBREW}
		
			% %
			% \subparagraph{\texorpdfstring{\he{sample word}}{}} % paste actual hebrew text here
			% %
				% \label{PAGA} % whatever is in step bible without periods
			
				% \begin{compactdesc}
					% \item[HALOT , PDF ] \textbf{} % Use the 1994 PDF
						% \begin{compactitem}
							% \item[1]
						% \end{compactitem}
					% \item[BDB , PDF ] \textbf{}
						% \begin{compactitem}
							% \item[1]
						% \end{compactitem}
					% \item[DCH PDF ] \textbf{} % don't include the actual page number, they repeat from volume to volume
						% \begin{compactitem}
							% \item[1]
						% \end{compactitem}
				% \end{compactdesc}
				
				% \begin{compactdesc}
					% \item[Some OT uses] \textbf{}
						% \begin{compactdesc}
							% \item[]
						% \end{compactdesc}
				% \end{compactdesc}
			
		% % --
		% \paragraph{Septuagint}
		% % --
			% \label{LOVE-LXX}
		
			% %
			% \subparagraph{\texorpdfstring{\gr{sample word}}{}}
			% %
		
	% % -----
	% \subsubsection{New Testament} % *NT
	% % -----
		% \label{LOVE-NT}
	
		% % --
		% \paragraph{KJV}
		% % --
			% \label{LOVE-NT-KJV}		
	
			% \begin{tabular}{ll}
				% Total occurrences in the NT & 
			% \end{tabular}
			
			% \begin{compactdesc}
				% \item[]
			% \end{compactdesc}
			
		% % --
		% \paragraph{Greek New Testament} % *GREEK
		% % --
			% \label{LOVE-GREEK}
		
			% %
			% \subparagraph{\texorpdfstring{\gr{sample word}}{}}
			% %
				% \label{KATALLAGE}
			
				% \begin{compactdesc}
					% \item[BDAG ] \textbf{}
						% \begin{compactitem}
							% \item 
						% \end{compactitem}
					% \item[CGL , PDF ] \textbf{}
						% \begin{compactitem}
							% \item 
						% \end{compactitem}
					% \item[LSJ , PDF ] \textbf{}
						% \begin{compactitem}
							% \item 
						% \end{compactitem}
				% \end{compactdesc}
				
				% \begin{tabular}{ll}
					% Total occurrences in the NT & 
				% \end{tabular}
				
				% \begin{compactdesc}
					% \item[Some NT uses] \textbf{}
						% \begin{compactdesc}
							% \item[]
						% \end{compactdesc}
				% \end{compactdesc}
				
			% %
			% \subparagraph{TDNT: \texorpdfstring{\gr{sample word}}{}  (3:536--556)}
			% %
				% \label{KATALLAGE-TDNT}
		
	% % -----
	% \subsubsection{Book of Mormon} % *BOM
	% % -----
		% \label{LOVE-BOM}
	
		% \begin{tabular}{ll}
			% Total occurrences in the BOM & 
		% \end{tabular}
		
		% \begin{compactdesc}
			% \item[]
		% \end{compactdesc}
		
	% -----
	\subsubsection{Doctrine and Covenants} % *DC
	% -----
		\label{LOVE-DC}
	
		% \begin{tabular}{ll}
			% Total occurrences in the DC & 
		% \end{tabular}
		
		\begin{compactdesc}
			\item[{\hyperref[DC95-12]{DC 95:12}}] If you keep not my commandments, the love of the Father shall not continue with you, therefore you shall walk in darkness.
				\begin{compactitem}
					\item This is a verse that you could use to talk about the idea of unconditional love, or the lack of it, as I think President Nelson did years ago when he was an apostle. See my notes there.
				\end{compactitem}
		\end{compactdesc}
		
	% % -----
	% \subsubsection{Pearl of Great Price} % *PGP
	% % -----
		% \label{LOVE-PGP}
	
		% \begin{tabular}{ll}
			% Total occurrences in the PGP & 
		% \end{tabular}
		
		% \begin{compactdesc}
			% \item[]
		% \end{compactdesc}
		
% % ----------------------
% \subsection{Key Phrases} % *KEYPHRASES *PHRASES
% % ----------------------
	% \label{LOVE-PHRASES}

	% \begin{compactdesc}
		% \item[] \textbf{\#times} | list
			% % \begin{compactdesc}
				% % \item[Bible anchor verses] \phantom{.}
					% % \begin{compactdesc}
						% % \item[Romans 5:5] \gr{}
					% % \end{compactdesc}
				% % \item[Commentary] ---
			% % \end{compactdesc}
	% \end{compactdesc}
	
% % ----------------------
% \subsection{Bible Dictionaries} % *DICTIONARIES
% % ----------------------
	% \label{LOVE-BD}

% % ----------------------
% \subsection{Restoration Usage} % *RESTORATION
% % ----------------------
	% \label{LOVE-RESTORATION}

	% % -----
	% \subsubsection{Joseph Smith} % *JS
	% % -----
		% \label{LOVE-JS}
	
		% \begin{compactdesc}
			% \item[{\hyperref[KEY]{Date/Source}}]
		% \end{compactdesc}
		
	% % -----
	% \subsubsection{Temple Ordinances}
	% % -----
		% \label{LOVE-TEMPLE}
	
		% \begin{compactdesc}
			% \item[{\hyperref[KEY]{Date/Source}}]
		% \end{compactdesc}
	
	% % -----
	% \subsubsection{Brigham Young} % *BY
	% % -----
		% \label{LOVE-BY}
	
		% \begin{compactdesc}
			% \item[{\hyperref[KEY]{Date/Source}}]
		% \end{compactdesc}
	
% ----------------------
\subsection{Commentary and Studies} % *COMMENTARY *STUDIES
% ----------------------
	\label{LOVE-COMMENTARY}

	% % -----
	% \subsubsection{Key Points}
	% % -----
		% \label{LOVE-KEYS}
	
		% \begin{compactitem}
			% \item
		% \end{compactitem}
		
	% -----
	\subsubsection{On the commandment to ``love one another''}
	% -----
		\label{LOVE-another}
		
		% --
		\paragraph{Augustine, DDC I:40--43}
		% --
		
			(See also I:67--69.)
		
			It is therefore an important question, whether humans should enjoy one another or use one another, or both. We have been commanded to love one another, but the question is whether one person should be loved by another on his own account or for some other reason. If on his own account, we enjoy him; if for another reason, we use him. In my opinion, he should be loved for another reason. For if something is to be loved on its own account, it is made to constitute the happy life, even if it is not as yet the reality but the hope of it which consoles us at this time. But `cursed is he who puts his hope in a man'. 41. Neither should a person enjoy himself, if you think closely about this, because he should not love himself on his own account, but only on account of the one who is to be enjoyed. A person is at his best when in his whole life he strives towards the unchangeable form of life and holds fast to it wholeheartedly. But if he loves himself on his own account, he does not relate himself to God, but turns to himself and not to something unchangeable. And for this reason it is with a certain insufficiency that he enjoys himself, because when totally absorbed and controlled by the unchangeable good he is a better man than when his attention leaves it, even if it turns to himself. 42. So if you ought to love yourself not on your own account but on account of the one who is the most proper object of your love, another person should not be angry if you love him too on account of God. For the divinely established rule of love says `you shall love your neighbour as yourself' but God `with all your heart, and with all your soul, and with all your mind', so that you may devote all your thoughts and all your life and all your understanding to the one from whom you actually receive the things that you devote to him. 43. And when it says `all your heart, all your soul, all your mind', it leaves no part of our life free from this obligation, no part free as it were to back out and enjoy some other thing; any other object of love that enters the mind should be swept towards the same goal as that to which the whole flood of our love is directed. So a person who loves his neighbour properly should, in concert with him, aim to love God with all his heart, all his soul, and all his mind. In this way, loving him as he would himself, he relates his love of himself and his neighbour entirely to the love of God, which allows not the slightest trickle to flow away from it and thereby diminish it. 